% Generated by scripts/assemble.py; corrections belong in source/corrections.json.
\examhead{8}{北京工业大学 2024 年数学分析试卷}{14}
% p015-b15; source PDF page 15
\par\Needspace{5\baselineskip}
\blockanchor{p015-b15}
\majorquestion{一}用致密性定理(有界数列必有收敛子列)证明闭区间连续函数可以达到最大值.\par
% p015-b16; source PDF page 15
\par\Needspace{5\baselineskip}
\blockanchor{p015-b16}
\majorquestion{二}$z = (x,y)$ 是方程 $x + y + z = \mathrm{e}^z$ 确定的隐函数，求 $\frac{\partial^2z}{\partial x\partial y}$.\par
% p015-b17; source PDF page 15
\par\Needspace{5\baselineskip}
\par\Needspace{6\baselineskip}
\blockanchor{p015-b17}
\majorquestion{三}求曲线积分\par
\nopagebreak[4]
% p015-b18; source PDF page 15
\blockanchor{p015-b18}
\displayformula{I = \int_ {L} \frac {y \mathrm{d} x - x \mathrm{d} y}{x ^ {2} + y ^ {2}}.}
% p015-b19; source PDF page 15
\blockanchor{p015-b19}
其中 L 是  $\frac{x^{2}}{4} + \frac{y^{2}}{9} = 1$  从 (2,0) 逆时针到 (-2,0) 的一段曲线.\par
% p015-b20; source PDF page 15
\par\Needspace{5\baselineskip}
\par\Needspace{6\baselineskip}
\blockanchor{p015-b20}
\majorquestion{四}用定积分求极限\par
\nopagebreak[4]
% p015-b21; source PDF page 15
\blockanchor{p015-b21}
\displayformula{\lim _ {n \to \infty} \sum_ {k = 1} ^ {n} \frac {n}{n ^ {2} + k} \sin \frac {k}{n}.}
% p015-b22; source PDF page 15
\par\Needspace{5\baselineskip}
\par\Needspace{6\baselineskip}
\blockanchor{p015-b22}
\majorquestion{五}已知在 $f(x, y)$ 在区域\par
\nopagebreak[4]
% p015-b23; source PDF page 15
\blockanchor{p015-b23}
\displayformula{D = \{(x, y) \mid a \leq x <   + \infty , c \leq y \leq d \}}
% p015-b24; source PDF page 15
\par\Needspace{6\baselineskip}
\blockanchor{p015-b24}
内连续, 且无穷积分 $\varphi(y) = \int_{a}^{+\infty} f(x, y) \, \mathrm{d}x$ 在 $[c, d]$ 上一致收敛, 试证明 $\varphi(y)$ 在 $[c, d]$ 上连续.
\majorquestion{六}计算\par
\nopagebreak[4]
% p015-b25; source PDF page 15
\blockanchor{p015-b25}
\displayformula{I = \int_ {0} ^ {+ \infty} \frac {\mathrm{e} ^ {- x} \sin a x}{x} \mathrm{d} x.}
% p016-b00; source PDF page 16
\par\Needspace{5\baselineskip}
\blockanchor{p016-b00}
\majorquestion{七}证明无穷积分 $\int_{1}^{+\infty} \frac{\sin x}{x^p} \cos \frac{1}{x} \mathrm{d}x$ 在 $p > 1$ 时绝对收敛, $0 < p \leq 1$ 时条件收敛.\par
% p016-b01; source PDF page 16
\par\Needspace{5\baselineskip}
\par\Needspace{6\baselineskip}
\blockanchor{p016-b01}
\majorquestion{八}计算\par
\nopagebreak[4]
% p016-b02; source PDF page 16
\blockanchor{p016-b02}
\displayformula{\oint_ {C} \left(y ^ {2} - z ^ {2}\right) \mathrm{d} x + \left(z ^ {2} - x ^ {2}\right) \mathrm{d} y + \left(x ^ {2} - y ^ {2}\right) \mathrm{d} z,}
% p016-b03; source PDF page 16
\blockanchor{p016-b03}
其中 $C: \left\{ \begin{array}{l} x^2 + y^2 + z^2 = 1 \\ x + y + z = 1 \end{array} \right.$ 的方向和 $x + y + z = 1$ 的法方向 $\vec{n} = (1, 1, 1)$ 构成右手螺旋系.\par
% p016-b04; source PDF page 16
\par\Needspace{5\baselineskip}
\par\Needspace{6\baselineskip}
\blockanchor{p016-b04}
\majorquestion{九}计算\par
\nopagebreak[4]
% p016-b05; source PDF page 16
\blockanchor{p016-b05}
\displayformula{I = \int_ {0} ^ {\frac {1}{2}} \mathrm{d} x \int_ {x} ^ {1 - x} \sqrt {y ^ {2} - x ^ {2}} \mathrm{d} y.}
% p016-b06; source PDF page 16
\par\Needspace{5\baselineskip}
\par\Needspace{6\baselineskip}
\blockanchor{p016-b06}
\majorquestion{十}已知 $f(x)$ 在 $[0,1]$ 存在连续导函数, $d_n = \int_0^1 f(x) \, \mathrm{d}x - \frac{1}{n} \sum_{k=1}^{n} f\left(\frac{k-1}{n}\right)$, 证明:\par
\nopagebreak[4]
% p016-b07; source PDF page 16
\blockanchor{p016-b07}
\displayformula{\lim _ {n \rightarrow \infty} n d _ {n} = \frac {1}{2} [ f (1) - f (0) ].}
